\section*{Summary}

A hot carrier in a semiconductor is a projectile of the master formula (\cref{sec:18.1}). Its impact ionization is the rate \cref{eq:18.1-ii}, whose kernel is the screened interaction, and phonon emission is the same formula with the electron--phonon vertex; momentum conservation puts the impact-ionization threshold at $\theta_{\mathrm{ii}}\Eg$, with $\theta_{\mathrm{ii}} = 3/2$ for equal masses. The cascade is a branching process whose statistics follow from one generating functional and its backward equation \cref{eq:18.1-backward}, and the energy identity \cref{eq:18.1-energy} ties every fluctuation of the pair number to a fluctuation of the heat. The cascade must be seeded class by class: at \SI{10}{eV} an equal split of the energy between electron and hole yields 2.985 pairs on average, and a uniform distribution of splits with the same mean split yields 2.691.

The mean energy per pair and the Fano factor are the first two moments of that generating functional (\cref{sec:18.2}). $F_{\mathrm{eh}}$ is the variance of the heat per pair in units of $\Eg^2$; in silicon its three terms add to $0.092 + 0.290 - 0.268 = 0.114$. Below about $15\Eg$ the yield rises in steps of one pair, so that a search at a few pairs needs the whole pair-number distribution.

At \SI{5}{eV} in silicon a dark-matter kick at $q_0$ puts two pairs in \SI{17}{\percent} of events, where the step model puts two in every event (\cref{sec:18.3}). The step model and fixed splits of the energy fail between 4.4 and \SI{13}{eV}, where the split decides the count. Against a full-band cascade with no fitted constant, the step model overcounts the two-pair bin about tenfold, and at least sixfold anywhere in the envelope of the full-band variants.

With its one constant fixed by the measured pair energy of silicon, the cascade yields $F_{\mathrm{eh}} = 0.114$ against the measured $0.119 \pm 0.002$ (\cref{sec:18.4}). The full-band cascade gives $\varepsilon_{\mathrm{eh}} = \SI{4.15}{eV}$ $[3.95, 4.73]$ against the measured \SI{3.749}{eV}. The calibrated cascade reduces over its parameter range to the observed linear law $\varepsilon_{\mathrm{eh}} \approx 1.61\Eg + k_2$. Against the seven crystals besides silicon with a measured pair energy, the residuals of the calibrated cascade run from $-\SI{19.8}{\percent}$ to $+\SI{7.0}{\percent}$ (root mean square \SI{10.5}{\percent}). The crystals with a measured pair energy, silicon included, span the validated range of the calibrated cascade, and the halides, fluorides, lead salts and InSb lie outside that range.

Below the impact-ionization threshold the pair number can still change (\cref{sec:18.5}). A capture probability $f_{\mathrm{cap}}$ shared by all electrons turns $\varepsilon_{\mathrm{eh}}$ into $\varepsilon_{\mathrm{eh}}/(1 - f_{\mathrm{cap}})$ and $F_{\mathrm{eh}}$ into $(1 - f_{\mathrm{cap}})F_{\mathrm{eh}} + f_{\mathrm{cap}}$, and a loss of \SI{0.54}{\percent} would account for the difference between the cascade value and the measured Fano factor of silicon. Excitons change neither the pair number nor its moments; they return their binding energy as heat, and an exciton that is never split thins the measured charge, \cref{eq:18.5-Fmeas}. Almost none of the rates of capture, exciton binding and trapping is known at millikelvin temperatures.
